\section{Introduction}

Why do HuggingFace datasets show 61\% preprocessing code documentation while UCI datasets show 0\% --- when both host ML datasets for the same research community? HuggingFace datasets show 55-66\% presence for preprocessing code and data provenance fields, while UCI ML Repository datasets show 0-15\% presence for the same fields --- a 45-61 percentage point gap that directly undermines reproducibility and reuse. A researcher attempting to reproduce a UCI dataset's preprocessing workflow finds only ``methodology text'' with no executable code, no dependency specifications, and no versioning metadata --- making automated reconstruction impossible.

Dataset reproducibility depends on documentation completeness, yet prior work treats incomplete metadata as creator failure rather than examining whether repository design systematically shapes documentation outcomes. Yang et al.~\cite{yang2024dataset} analyzed HuggingFace dataset cards at scale (7,433 datasets), revealing marked heterogeneity: Dataset Description and Structure sections show higher completion than Considerations sections, with completion correlating with dataset popularity. Strecker~\cite{strecker2026metadata} identified metadata conflicts across 8 geoscience and social science repositories, attributing incomplete DataCite metadata to both implementation workflows and inter-standard schema differences. However, no study has quantitatively linked repository UX features to metadata completeness patterns across platforms.

We introduce \textbf{friction score} (0-4 scale: automated field extraction + pre-filled templates + validation feedback + programmatic API access) as a quantifiable UX metric predicting metadata presence patterns. Our key insight is that repository friction-reduction features enable voluntary completion for optional metadata fields (45-61pp effect sizes), while enforcement mechanisms maintain required field presence ($\sim$75-95\%) --- two magnitude-differentiated pathways serving distinct completeness goals. Platforms with higher friction scores (HuggingFace friction=3: templates + validation + API) show 55-66\% optional field presence versus UCI (friction=0: manual web forms only) showing 0-15\%, with effect sizes of 45-61 percentage points. Required fields (license, version) show 75-95\% presence across ALL platforms with weaker friction effects (12-20pp, Cramér's $V$ 0.1-0.2), validating enforcement as a distinct mechanism.

Building on this insight, we make the following contributions:

\textbf{First cross-platform friction study:} We extend Yang (2024)'s single-platform analysis to 3 ML repositories (OpenML, HuggingFace, UCI) at 10,000+ dataset scale, introducing friction score as explanatory variable and testing the specific mechanism distinguishing enforcement from friction reduction.

\textbf{Mechanism validation:} We demonstrate that enforcement sets floor (UCI 74-75\% social norm, HuggingFace/OpenML 89-95\% technical blocking) while friction reduction raises ceiling (90\%$\rightarrow$95\%). Within-platform comparison (HuggingFace API vs manual uploads) confirms friction features causally reduce entry cost (12.1pp higher completeness, $p<0.0001$, Cohen's $d=0.621$).

\textbf{Quantitative effect sizes:} We measure 45-61pp differences for optional fields and 12-20pp for required fields, showing practical significance beyond statistical significance. Gradient effects validated for 2/3 optional fields (data\_source\_url, collection\_date show UCI $<$ OpenML $<$ HuggingFace), with threshold effect for code fields (preprocessing\_code requires platform code snippet support).

\textbf{Repository design insights:} We provide actionable recommendations from gradient analysis: combine enforcement (required fields) with friction-reduction UX (optional fields). Estimated marginal effects suggest templates (est.\ 25-30pp) $>$ API (est.\ 10-15pp) $>$ validation (est.\ 5-10pp), though feature ablation experiments are needed to validate these decomposed effects.

The remainder of this paper is organized as follows: Section~\ref{sec:related} reviews related work on dataset documentation and metadata standardization. Section~\ref{sec:methodology} describes our friction scoring methodology and large-scale extraction approach. Section~\ref{sec:experiments} details experimental design across 4 sub-hypotheses. Section~\ref{sec:results} presents results validating friction-completeness correlation. Section~\ref{sec:discussion} discusses mechanism interpretation, honest limitations, and broader impact. Section~\ref{sec:conclusion} concludes with future directions.
